\documentclass[../../../include/open-logic-section]{subfiles}

\begin{document}

\olfileid{inc}{art}{int}
\olsection{引言}

% Part: incompleteness
% Chapter: arithmetization-syntax
% Section: introduction

为了将可计算性与逻辑联系起来，我们需要一种便于计算处理的方式来谈论逻辑对象（符号、项、公式、推导），以及关于它们的操作、性质和关系。
我们可以直接实现这一点，即考虑符号、符号序列以及其他由它们构造出的对象上的可计算函数和关系。
由于逻辑语法的对象都是有限的，并且是从可数的符号集合构造出来的，这对于某些计算模型是可行的。
但其他计算模型——例如递归函数——则仅限于处理数、数的关系和函数。
此外，我们最终还希望能处理某些理论内部的语法，具体来说，是那些以算术语言表述的理论。
在这些情况下，就必须对语法进行\emph{算术化}，也就是说，分别将语法对象、对它们的操作以及它们之间的关系表示为数字、算术函数和算术关系。
这一思想可追溯到 莱布尼茨，其核心理念就是给语法对象分配数字。

给符号分配数字作为其“代码”相对简单。
有些符号会带来一点挑战，因为，例如，存在无穷多个变元，甚至对于每个元数~$n$ 都有无穷多个函数符号。
但当然，我们能够系统地给符号分配数字，使得，例如，$\Obj v_2$ 和 $\Obj v_3$ 被分配不同的代码。
符号序列（如项和公式）则是更大的挑战。
但如果我们能纯粹用算术方式处理数的序列（例如，通过序列的素数幂编码），那么我们就可以把对单个符号的编码扩展到对符号序列的编码，并进一步扩展到公式的序列或其他排列方式，如推导。这种扩展的编码被称为“哥德尔编码”。
每个项、公式和推导都被分配一个哥德尔数。

通过将符号序列编码为其代码的序列，并选用一种能用可计算函数处理的序列编码系统，我们就可以用可计算函数来处理哥德尔数了。
实际上，所有相关的函数都将是原始递归的。
例如，从序列的代码计算序列的长度，以及计算序列的第 $i$ 个元素，都是原始递归的。
如果编码某个序列的数恰好是公式~$!A$ 的哥德尔数，我们立刻就能看出，公式的长度和公式中第 $i$ 个符号的（代码）也都可以从~$!A$ 的哥德尔数计算出来。
要证明某些性质——例如，是一个合法构成的项的哥德尔数，或是一个正确推导的哥德尔数——是原始递归的，则要稍难一些。
但这仍然是可能的，因为我们感兴趣的序列（项、公式、推导）都是归纳定义的。

举个例子，考虑代入操作。
如果 $!A$ 是一个公式，$x$ 是一个变元，$t$ 是一个项，那么 $\Subst{!A}{t}{x}$ 就是将~$!A$ 中 $x$ 的所有自由出现替换为~$t$ 的结果。
现在假设我们给 $!A$, $x$, $t$ 分配了哥德尔数——设分别为 $k$, $l$ 和 $m$。
同一个编码方案会给 $\Subst{!A}{t}{x}$ 分配一个哥德尔数，设为~$n$。
这个将 $k$, $l$, $m$ 映射到 $n$ 的映射，就是代入操作在算术上的对应物。
当代入操作将 $!A$, $x$, $t$ 映射到 $\Subst{!A}{t}{x}$ 时，算术化的代入函数就将哥德尔数 $k$, $l$, $m$ 映射到哥德尔数~$n$。
我们将会看到，这个函数是原始递归的。

语法的算术化并非只有抽象的理论趣味。
尽管最初它是一个非凡的洞见，揭示了像算术语言这样本身不具备“谈论”语言机制的语言，毕竟也能形式化表达式的复杂性质。
由此自然引出一个问题：用这种语言表述的理论，例如 皮亚诺算术，能对其自身语言\emph{证明}些什么（例如，某个语句是否可证或是否为真）。
这就将我们引向了著名的限制性定理：哥德尔的（关于不可证性的）定理和 Tarski（塔斯基）的（真之不可定义性）定理。
不仅如此，算术化语法这一技巧对于证明可计算性理论中的一些重要结果也至关重要，例如关于理论的计算能力，或不同计算模型之间的关系。
语法的算术化也为算术化其他对象和性质提供了一个范例。
例如，我们同样可以将（比方说，图灵机的）格局与计算算术化。
这使得在一个计算模型（如图灵机）中的计算能在另一个计算模型（如递归函数）中被模拟。

\end{document}